\documentclass[../../../ps_q.tex]{subfiles}

\begin{document}
    \setcounter{figure}{0}
    図1のように，真空中に，極板面積$S$，極板間隔$d$の平行平板コンデンサがあり，%
    起電力$V_0$の電池をつないである．真空の誘電率を$\varepsilon_0$とし，端の効果は無視する．

    \begin{enumerate}
        \item 図\ref{fig:fig_em_cx_10_01_q}のように，コンデンサの極板間のちょうど下半分の領域を%
            比誘電率$\varepsilon_r$の誘電体で満たし，安定するまで待った．%
            このとき，コンデンサに蓄えられた電荷の大きさ$Q_1$，および%
            真空部品の電場の強さ$E_1$を求めよ．
        \item 図\ref{fig:fig_em_cx_10_02_q}のように，コンデンサの極板間のちょうど右半分の領域を%
            比誘電率$\varepsilon_r$の誘電体で満たし，安定するまで待った．%
            このとき，コンデンサに蓄えられた電荷の大きさ$Q_2$，および%
            真空部品の電場の強さ$E_2$を求めよ．
    \end{enumerate}
    \vspace{10pt}
    \begin{minipage}{\linewidth}
        \setlength{\abovecaptionskip}{2pt}
        \centering
            \includegraphics{\subfix{fig_em_cx_10/fig_em_cx_10_01_q.pdf}}
            \captionof{figure}{} \label{fig:fig_em_cx_10_01_q}
            \vspace{12pt}
            \includegraphics{\subfix{fig_em_cx_10/fig_em_cx_10_02_q.pdf}}
            \captionof{figure}{} \label{fig:fig_em_cx_10_02_q}
            \vspace{12pt}
    \end{minipage}
\end{document}






